\par\Needspace{4\baselineskip}
\begin{center}
\SourceLeaf{029}\pdfbookmark[2]{11.}{article-011}11.
\end{center}
Das inverse Problem, das unter dem Namen \textsc{Kepler}\emph{sches Problem} bekannt ist und darin besteht, aus der mittleren Anomalie die wahre Anomalie und den Radiusvektor zu bestimmen, kommt in der Praxis viel häufiger vor. Die Astronomen stellen die Mittelpunktsgleichung gewöhnlich als unendliche Reihe in den Sinuswerten der Winkel $M$, $2M$, $3M$ usw. dar. Jeder Koeffizient eines solchen Sinusterms ist selbst eine unendliche Reihe in Potenzen der Exzentrizität. Mehrere Autoren haben diese Formel für die Mittelpunktsgleichung entwickelt. Wir hielten es hier umso weniger für nötig, bei ihr zu verweilen, als sie sich unseres Erachtens für den praktischen Gebrauch weit weniger eignet als das indirekte Verfahren, besonders wenn die Exzentrizität nicht sehr klein ist. Daher werden wir das indirekte Verfahren in der Form, die uns am zweckmäßigsten erscheint, etwas ausführlicher erläutern.

Die Gleichung XII, $E = M + e\sin E$, gehört zu den transzendenten Gleichungen und lässt sich durch endlich viele direkte Rechenoperationen nicht lösen. Man muss sie durch wiederholtes Probieren lösen, indem man von einem Näherungswert für $E$ ausgeht und diesen mit geeigneten Verfahren so oft korrigiert, bis er der Gleichung genau genügt, also entweder mit der vollen Genauigkeit, die die Sinustafeln erlauben, oder wenigstens mit der für den beabsichtigten Zweck ausreichenden Genauigkeit. Erfolgen diese Korrekturen nach einer zuverlässigen und festen Regel, so besteht zwischen einem solchen indirekten Verfahren und einer Lösung durch Reihen kaum ein wesentlicher Unterschied. Lediglich der Anfangswert der Unbekannten ist beim indirekten Verfahren bis zu einem gewissen Grad frei wählbar. Das ist eher als Vorteil anzusehen, weil ein passend gewählter Wert die Korrekturen erheblich beschleunigen kann. Es sei $\epsilon$ ein Näherungswert für $E$ und $x$ die noch hinzuzufügende Korrektur (in Bogensekunden), sodass der Wert $E = \epsilon + x$ unserer Gleichung genau genügt. Man berechne $e\sin\epsilon$ in Bogensekunden mithilfe von Logarithmen und entnehme dabei den Tafeln zugleich die Änderung von $\log\sin\epsilon$ bei einem Schritt von $1^{\prime\prime}$ für den Winkel $\epsilon$ sowie die Änderung von $\log e\sin\epsilon$ bei einer Änderung der Zahl $e\sin\epsilon$ um eine Einheit. Diese Änderungen seien ohne Rücksicht auf ihre Vorzeichen der Reihe nach $\lambda$ und $\mu$. Es bedarf kaum der Erwähnung, dass beide Logarithmen mit gleich vielen Dezimalstellen angegeben sein sollen. Liegt nun $\epsilon$ bereits so nahe am wahren Wert von $E$, dass die Änderungen des Logarithmus des Sinus von $\epsilon$ bis $\epsilon+x$ und die Änderungen des Logarithmus der Zahl von $e\sin\epsilon$ bis $e\sin(\epsilon+x)$ als gleichmäßig gelten können, so kann man offenbar $e\sin(\epsilon+x) = e\sin\epsilon \pm \frac{\lambda x}{\mu}$ setzen, wobei das \SourceLeaf{030}obere Vorzeichen im ersten und vierten Quadranten und das untere im zweiten und dritten gilt. Da nun $\epsilon+x = M+e\sin(\epsilon+x)$ ist, ergibt sich $x = \frac{\mu}{\mu\mp\lambda}(M+e\sin\epsilon-\epsilon)$ und für den wahren Wert von $E$ folglich $\epsilon+x = M+e\sin\epsilon \pm \frac{\lambda}{\mu\mp\lambda}(M+e\sin\epsilon-\epsilon)$, mit den Vorzeichen wie angegeben. Ferner sieht man leicht, dass ohne Rücksicht auf das Vorzeichen $\mu:\lambda = 1:e\cos\epsilon$ und daher stets $\mu > \lambda$ gilt. Somit liegt im ersten und letzten Quadranten $M+e\sin\epsilon$ zwischen $\epsilon$ und $\epsilon+x$, im zweiten und dritten dagegen $\epsilon+x$ zwischen $\epsilon$ und $M+e\sin\epsilon$. Diese Regel kann es erleichtern, auf die Vorzeichen zu achten. War der angenommene Wert $\epsilon$ noch zu weit vom wahren Wert entfernt, um die obige Annahme als hinreichend genau gelten zu lassen, so ergibt dieses Verfahren gewiss einen wesentlich näheren Wert. Mit ihm ist dieselbe Rechnung noch einmal oder, falls nötig, mehrmals zu wiederholen. Betrachtet man den Unterschied zwischen dem Anfangswert $\epsilon$ und dem wahren Wert als Größe erster Ordnung, so ist leicht ersichtlich, dass der Fehler des neuen Werts von zweiter Ordnung ist und sich durch Wiederholung der Rechnung auf die vierte, achte und weitere Ordnungen vermindert. Je kleiner außerdem die Exzentrizität ist, desto schneller konvergieren die aufeinanderfolgenden Korrekturen.
